\documentclass[11pt]{article}
\usepackage[a4paper,margin=28mm]{geometry}
\usepackage{amsmath,amssymb,amsthm,booktabs,hyperref}
\newtheorem{theorem}{Theorem}[section]
\newtheorem{lemma}[theorem]{Lemma}
\newtheorem{corollary}[theorem]{Corollary}
\newtheorem{proposition}[theorem]{Proposition}
\newcommand{\dist}{\operatorname{dist}}
\title{A regular-graph bound at diameter eight and its degree-four equality case}
\author{Human Compression project: mathematical preparation}
\date{31 October 2026}
\begin{document}\maketitle

\paragraph{Scope and attribution.}
This is a written proof prepared from a hand argument proposed by the
project's coordinating agent. It uses the standard independent-set
degree bound and elementary graph distances. It does not use an original
literature census count, execute a graph enumeration, or certify an
original manuscript or a conjecture by computation or Lean. The
fixed-geodesic method is the common setup of the retained degree-five
normal forms; the argument here isolates the part that works for all
degrees. These new statements await an independent review and adoption.

All graphs in this note are finite, simple and connected. The girth
condition is triangle-freeness; the equality constructions will have
girth exactly four. A graph of diameter eight has a geodesic
$p_0,p_1,\ldots,p_8$ of length eight.

\section{The common bound}

\begin{lemma}\label{independent}
Let $\Gamma$ be a $k$-regular triangle-free graph of diameter eight,
where $k\geq2$, and let $n=|V\Gamma|$. Then $n\geq6k$.
If equality holds, the graph is bipartite with parts
\[
 I=N(p_0)\cup N(p_4)\cup N(p_8),\qquad J=V\Gamma\setminus I,
\]
each of size $3k$.
\end{lemma}
\begin{proof}
Put $A=N(p_0)$, $B=N(p_4)$ and $C=N(p_8)$. Each set is independent
by triangle-freeness. The three sets are disjoint: a vertex in two of
them would give a path of length two between the corresponding
geodesic vertices, whose mutual distances are four or eight. There
are no edges between $A$ and $B$, or between $B$ and $C$, since such
an edge would give a path of length three between vertices at distance
four. An edge between $A$ and $C$ would give a path of length three
between $p_0$ and $p_8$. Thus $I$ is independent and $|I|=3k$.

All $3k^2$ edges incident with $I$ go into its complement. The latter
has total degree $k(n-3k)$, so $3k^2\leq k(n-3k)$ and $n\geq6k$.
If $n=6k$, every degree in the complement is used by an edge to $I$;
there is no edge within $J$. This gives the stated bipartition.
\end{proof}

Assume for the moment that $n=6k$. The three vertices $p_0,p_4,p_8$
lie in $J$ and are distinct. Define
\[
 S_0=\{x\in J\setminus\{p_0,p_4,p_8\}:N(x)\cap A\ne\varnothing\},
 \quad
 S_8=\{x\in J\setminus\{p_0,p_4,p_8\}:N(x)\cap C\ne\varnothing\}.
\]
They are disjoint: a common vertex would give a $p_0$--$p_8$ path
of length four. Let $M$ be the remaining vertices of
$J\setminus\{p_0,p_4,p_8\}$, and put $m=|M|$.

\begin{lemma}\label{balance}
In this equality setup, write $s_i=|S_i|$ for $i=0,8$. Then
\[
 s_0,s_8\geq k,\qquad 0\leq m\leq k-3.
\]
Every member of $M$ is adjacent to every member of $B$. Put
$q=k-1-m$, so $2\leq q\leq k-1$. The set $B$ splits into nonempty
sets $B_0,B_8$, with each vertex in $B_i$ having exactly $q$ neighbors
in $S_i$ and none in the other side. If $b_i=|B_i|$, then
\begin{equation}\label{balanceeq}
 b_0+b_8=k,\qquad q b_i=k(s_i-k+1)\quad(i=0,8).
\end{equation}
\end{lemma}
\begin{proof}
Each vertex of $A$ has $p_0$ as one neighbor and its other $k-1$
neighbors in $S_0$. Indeed it cannot meet $p_4,p_8$ by the geodesic
distances, and the definition includes every other possible neighbor.
Hence $s_0\geq k-1$. If equality held, all vertices of $A$ would be
adjacent to all of $S_0$, and each member of $S_0$ would already have
its full $k$ neighbors in $A$. The vertices $A\cup S_0\cup\{p_0\}$
would form a connected component $K_{k,k}$, contradicting connectedness
and the existence of the distant endpoint. Thus $s_0\geq k$. The
same proof applies on the other side. Since
$s_0+s_8+m=3k-3$, it follows that $m\leq k-3$.

A vertex of $M$ has no neighbor in $A$ or $C$, by definition, and all
its neighbors lie in $I$. It therefore has exactly the $k$ vertices
of $B$ as neighbors. Each vertex of $B$ is consequently adjacent to
$p_4$ and all $m$ vertices of $M$, leaving exactly $q=k-1-m$ neighbors
in the sides. This number is positive, and the bounds on $m$ give
$2\leq q\leq k-1$. A vertex of $B$ cannot meet both sides: otherwise
there is a path $p_0,A,S_0,B,S_8,C,p_8$ of length six. Thus it belongs
to a well-defined $B_i$.

The total number of edges from $A$ to $S_0$ is $k(k-1)$, while all
degrees in $S_0$ total $ks_0$. Their difference must be supplied by
$B_0$, so $qb_0=ks_0-k(k-1)$. This is positive since $s_0\geq k$,
proving also $b_0>0$. The other side is identical. All vertices of
$B$ have been assigned, so $b_0+b_8=k$.
\end{proof}

\begin{theorem}\label{prime}
Let $\Gamma$ be a triangle-free $k$-regular graph of diameter eight.
If $k$ is prime, then $|V\Gamma|>6k$. If $k$ is an odd prime, then
\[
 |V\Gamma|\geq6k+2.
\]
In particular every quintic triangle-free graph of diameter eight
has at least $32$ vertices.
\end{theorem}
\begin{proof}
Lemma~\ref{independent} gives the bound $6k$. For $k=2$, equality
would contradict $0\leq m\leq-1$ in the proof of Lemma~\ref{balance}.
For a prime $k\geq3$, equality would give the balances
\eqref{balanceeq}. Since $2\leq q\leq k-1$, the prime $k$ does not
divide $q$. It would therefore divide each $b_i$, whereas
$1\leq b_i\leq k-1$. This contradiction excludes equality.
For odd $k$, the degree sum $kn$ is even, so $n$ is even. The next
possible value after $6k$ is at least $6k+2$.
\end{proof}

This theorem proves only the lower bound. To conclude that the quintic
minimum equals $32$, one must separately import or provide one actual
quintic graph of order $32$, girth four and diameter eight. An existing
reviewed construction may supply that upper bound; no literature count
is used in this proof. The lower orders $26,28,30$ need no separate
generation or exclusion computations once this common theorem is used.

\section{The degree-four equality data}

\begin{lemma}\label{fourdata}
If a triangle-free quartic graph has diameter eight and order $24$,
then in the preceding setup
\[
 |S_0|=|S_8|=4,\quad m=1,\quad q=2,\quad |B_0|=|B_8|=2.
\]
The two vertices in $X=\{p_4\}\cup M$ are false twins, with common
neighborhood $B$. Removing $X$ leaves two connected components, each
of order eleven. In each component exactly two vertices have degree
two, namely its vertices of $B_i$; every other vertex has degree four.
\end{lemma}
\begin{proof}
Lemma~\ref{balance} gives $m\leq1$. If $m=0$, then $q=3$, and
$3b_i=4(s_i-3)$ forces $4\mid b_i$, impossible for
$1\leq b_i\leq3$. Thus $m=1$ and $q=2$. The balance now forces
$b_i$ even, hence both are two, and then $s_i=4$.

The single vertex of $M$ and $p_4$ both have neighborhood $B$. There
are no edges between the two sides, and no edges within $J$. Each
vertex of $A$ has three neighbors in $S_0$. A side vertex has degree
four and at most two neighbors in $B_0$, hence at least two in $A$.
Thus the endpoint $p_0$, the four vertices of $A$, and all of $S_0$
lie in one connected component after removing $X$; the two ports in
$B_0$ each have two side neighbors and lie in this component as well.
It has $1+4+4+2=11$ vertices. The other side is identical. A port
loses its two neighbors in $X$, so has degree two in its component.
No other vertex has a neighbor in $X$, so all other internal degrees
remain four.
\end{proof}

\paragraph{An eleven-vertex half.}
Let $W$ be a four-element set. Choose two two-element subsets
$P,Q\subseteq W$, permitting $P=Q$. Create two distinct ports $u,v$
with neighborhoods $P,Q$. Create an endpoint $e$ and four row vertices
$r_1,\ldots,r_4$, joining $e$ to all rows. For $w\in W$, put
\[
 \beta(w)=\boldsymbol1_P(w)+\boldsymbol1_Q(w).
\]
Choose the multiset of omitted side vertices of the four rows so that
$w$ is omitted exactly $\beta(w)$ times; this is possible because
$\sum_w\beta(w)=4$. Join each row to the other three side vertices.
Rows with equal omitted vertices are still distinct vertices. There
are no other edges. Denote this graph, with its intrinsic two-port
set, by $H(P,Q)$.

\begin{lemma}\label{half}
The graphs $H(P,Q)$ are exactly the connected components in
Lemma~\ref{fourdata}, as unrooted graphs. They give exactly three
ordinary isomorphism classes, parametrized by
$t=|P\cap Q|\in\{0,1,2\}$.
\end{lemma}
\begin{proof}
In a component of Lemma~\ref{fourdata}, let $P,Q$ be the two ports'
side neighborhoods. Each has size two. Each of the four rows has
three side neighbors and therefore omits exactly one side vertex.
The degree of a side vertex requires its row incidence to be
$4-\beta(w)$; hence its omission multiplicity is exactly $\beta(w)$.
This proves necessity and also shows that no further row choices
remain, apart from permuting the row vertices.

Conversely the construction gives each side vertex degree
$(4-\beta(w))+\beta(w)=4$, every row degree $1+3=4$, the endpoint
degree four and each port degree two. Every side vertex meets at
least two rows, so the graph is connected. It is bipartite, with
parts consisting of rows and ports on one side and $e\cup W$ on
the other.

Under permutations of $W$ and exchange of the two ports, two pairs
$(P,Q)$ are equivalent precisely when their intersection sizes are
equal: the sizes of $P\cap Q$, $P\setminus Q$, $Q\setminus P$ and
$W\setminus(P\cup Q)$ are respectively $t,2-t,2-t,t$.
An equivalence extends to the forced row multisets and gives an
isomorphism of the halves. Finally the ports are exactly the vertices
of degree two, so their number of common neighbors is an intrinsic
ordinary graph invariant. It equals $t$. Distinct values of $t$
therefore cannot be identified by an unrooted isomorphism, even one
which changes the displayed endpoint or row presentation.
\end{proof}

\section{Metric and ordinary gluing}

Take any two halves $H(P_0,Q_0)$ and $H(P_8,Q_8)$ on disjoint vertices.
Add two new nonadjacent vertices $x,x'$ and join both to each of the
four ports. Denote their set by $X$. There is no alignment parameter:
every port meets both new vertices.

\begin{lemma}\label{metric}
Every graph obtained by this gluing is simple, connected, quartic,
of girth four and diameter eight. Its center is exactly $X$.
\end{lemma}
\begin{proof}
Simplicity and connectedness follow from the construction. A port
of internal degree two gains its two neighbors in $X$, each member
of $X$ has four port neighbors, and all other degrees remain four.
The bipartitions of the halves combine with $X$ in the rows-and-ports
opposite part; the full graph is bipartite, so has no triangle.
Within either half the four row neighborhoods are three-subsets of
a four-set. Any two such neighborhoods have at least two common
vertices and form a four-cycle. Thus the girth is exactly four.

We compute distances to either member of $X$. Each port is at
distance one. A side vertex in $P\cup Q$ is at distance two.
Every row has a three-element side neighborhood, which meets
$P\cup Q$ since this union has size at least two. Hence every row
is at distance three. A side vertex outside $P\cup Q$ has
$\beta=0$ and is adjacent to all four rows, so is at distance four.
The endpoint $e$ is also at distance four. These distances have the
stated exact values by the layer structure and bipartiteness. In
particular all vertices are within four of $x$ and of $x'$; their
mutual distance is two. Every pair of graph vertices is consequently
at distance at most eight.

Removing $X$ leaves the two halves. The endpoint in either half
has distance four to $X$, so the two endpoints are at distance eight.
It follows that the diameter is eight and the eccentricities of
$x,x'$ are four.

Every other vertex $z$ lies in one half. The endpoint $e'$ of the
other half has distance four to both members of $X$. All paths
from $z$ to $e'$ pass through $X$. Since the two members of $X$
are false twins, their distances to $z$ agree. Hence
\[
 \dist(z,e')=\dist(z,X)+4>4.
\]
No such $z$ is central. Thus the graph center is exactly $X$.
\end{proof}

\begin{theorem}\label{six}
There are exactly six ordinary isomorphism classes of triangle-free
quartic graphs of order $24$ and diameter eight. Each has girth four.
They are the gluings of an unordered pair, with repetition allowed,
of the three half classes in Lemma~\ref{half}.
\end{theorem}
\begin{proof}
Lemmas~\ref{fourdata} and~\ref{half} prove that every graph in the
specified domain has this form. Lemma~\ref{metric} proves the
converse for every choice of halves.

An ordinary isomorphism preserves the graph center, since it preserves
all distances and eccentricities. Thus it preserves $X$ and sends the
two components after deleting $X$ to the two corresponding components,
possibly exchanging them. Their port sets are intrinsic by their
internal degrees. Therefore isomorphic full graphs have the same
unordered pair of half classes. Conversely component isomorphisms,
with a possible exchange of the two components and any exchange of
$x,x'$, extend to a full isomorphism: all ports meet all of $X$.
There are $\binom{3+2-1}{2}=6$ unordered pairs of three classes with
repetition allowed. No rooted marking or literature count enters this
ordinary-class calculation.
\end{proof}

\paragraph{What the proof replaces.}
The independent-set lemma charges one common geodesic argument for
all degrees. The equality balance and prime divisibility exclude all
orders at or below $6k$ for prime $k$. The degree-four equality case
requires only two small port neighborhoods and a forced omission
multiset. These are deductions from common statements, not claims
that every other result in a source manuscript follows from them.
No computational validation or Lean verification of this note has
been run.

\end{document}
